El 1877, l'astrònom Asaph Hall va descobrir els satèl·lits del planeta Mart: Fobos i Deimos. El dia 6 d'agost de 2012, el robot \textit{Curiosity} va arribar al planeta Mart i des de llavors envia informació a la Terra sobre les característiques d'aquest planeta. A partir de les dades subministrades, calculeu:

\figura[0.5]{fig1.png}

\begin{apartat}{1}
La massa del planeta Mart.
\end{apartat}

\begin{apartat}{1}
El radi de l'òrbita de Deimos i la velocitat d'escapament del robot \textit{Curiosity} des de la superfície del planeta.
\end{apartat}

\dades{%
  Radi de Mart, $R_\mathrm{Mart} = 3\,390\ \mathrm{km}$\\
  Acceleració de la gravetat en la superfície de Mart, $g_\mathrm{Mart} = 3,71\ \mathrm{m\,s^{-2}}$\\
  Període orbital de Deimos, $T_\mathrm{Deimos} = 30,35\ \mathrm{h}$\\
  Massa de Deimos, $m_\mathrm{Deimos} = 2\times 10^{15}\ \mathrm{kg}$\\
  $G = 6,67\times 10^{-11}\ \mathrm{N\,m^2\,kg^{-2}}$}
